\chapter{Methodology}
\label{ch:methodology}

\section{System overview}
\label{sec:overview}

The system is a five-layer stack (Figure~\ref{fig:architecture}) around one
governing idea: every claim an agent makes must trace back to something it
was given. The Python backend runs to 33 modules and roughly 4,900 lines
with a 92-case test suite; the API and frontend add 374 and 3,800 lines.

\begin{figure}[htbp]
\centering
\resizebox{\textwidth}{!}{%
\begin{tikzpicture}[font=\footnotesize,
 badge/.style={circle,draw=f21ink,fill=white,text=f21ink,
 font=\bfseries\footnotesize,minimum size=9mm,inner sep=0pt},
 layer/.style={draw=f21ink,line width=0.7pt,rounded corners=3pt,
 text width=10.8cm,minimum height=1.15cm,align=left,
 fill=f21mist,inner sep=7pt},
 gaplayer/.style={draw=f21grey,dashed,line width=0.7pt,rounded corners=3pt,
 text width=10.8cm,minimum height=0.95cm,align=left,
 fill=white,inner sep=7pt},
 layertag/.style={f21grey,font=\scriptsize\itshape}]

 \node[badge] (b1) at (0,0) {L1};
 \node[layer,right=6mm of b1,anchor=west] (l1)
 {\textbf{Presentation} --- React/TypeScript client: upload, live streamed\\
 negotiation, case browser, analytics, risk screen \hfill{\scriptsize\itshape\color{f21grey}browser}};

 \node[badge,below=8mm of b1] (b2) {L2.5};
 \node[layer,right=6mm of b2,anchor=west] (l2)
 {\textbf{API} --- FastAPI: extraction endpoint, server-sent-event\\
 negotiation stream, risk assessment, GPU-tunnel status and reconnect};

 \node[badge,below=8mm of b2] (b3) {L2};
 \node[layer,right=6mm of b3,anchor=west] (l3)
 {\textbf{Orchestration} --- LangGraph \texttt{StateGraph}: deterministic\\
 turn order, conditional routing, per-super-step streaming \hfill{\scriptsize\itshape\color{f21grey}process}};

 \node[badge,below=8mm of b3] (b4) {L3};
 \node[layer,right=6mm of b4,anchor=west] (l4)
 {\textbf{Agents} --- Contracting Authority, Aggrieved Bidder, Court,\\
 Summary, Extraction (Llama 3.1 8B via Ollama; local CPU or\\
 tunnelled A10G GPU) \hfill{\scriptsize\itshape\color{f21grey}model}};

 \node[badge,below=8mm of b4] (b5) {L4};
 \node[gaplayer,right=6mm of b5,anchor=west] (l5)
 {\textbf{Knowledge / RAG} --- vector retrieval over legislation and\\
 case law \emph{--- deliberately not built; see \S\ref{sec:alternatives}}};

 \node[badge,below=8mm of b5] (b6) {L5};
 \node[layer,right=6mm of b6,anchor=west] (l6)
 {\textbf{Persistence} --- structured JSON logs per run; batch summaries\\
 with provenance and compliance counters};

 \draw[->,f21ink] (b1) -- (b2);
 \draw[->,f21ink] (b2) -- (b3);
 \draw[->,f21ink] (b3) -- (b4);
 \draw[->,f21grey,dashed] (b4) -- (b5);
 \draw[->,f21ink] (b5) -- (b6);
\end{tikzpicture}}
\caption[System architecture]{The five-layer architecture, read top to bottom as
the path a request takes. The dashed layer is reported because it was scoped
out, not because it exists: the extraction agent is a single-pass synthesis
step and must not be mistaken for retrieval.}
\label{fig:architecture}
\end{figure}

\section{Grounding the agents in the client's own material}
\label{sec:grounding}

Assigning a model a role and interests is role play, not instantiation: it is
steered toward the region of its distribution that talks like a procurement
officer, and a persuasive performance is not evidence the reasoning is
sound. An early version used interest lists I had written myself, believable
but unchallengeable; these were replaced with the taxonomies in Fusion21's
Doc 1 (seven categories for the contracting authority, six each for the
supplier and judiciary), so every agent's stated interests trace to a named
category in a document the client wrote. Per-agent configuration ---
interest source, sampling temperature, structured output type --- is
recorded in Appendix~\ref{app:repro}; temperatures were set by role, held
low for the adjudicator and summariser where consistency matters more than
variety, and highest for the challenger, where a negotiation in which
nobody proposes anything new stalls.

Doc 1 also supplies per-case behavioural modifiers in its ``practical
considerations'' tables, varying by dispute rather than fixed to a role.
These are optional \texttt{CAProfile}/\texttt{BidderProfile} objects
attached to a scenario and rendered into the system prompt at construction;
both default to unset and render no text, so results obtained before
profiles existed stay comparable. Consequently agents can no longer be
module-level singletons, since their prompts depend on the scenario, so they
are built per dispute and bound into the graph nodes at compile time.

The fourth agent takes no part in the negotiation: it reads the finished
transcript and writes a plain-English account of what happened, on the
reasoning that a system only a procurement lawyer can parse has not
lowered costs. Section~\ref{sec:dynamics} tests that claim against a
standard readability measure --- the summaries read fluently enough that
the shortfall is invisible without measuring it.

\section{The Court agent: conditional verification}
\label{sec:court-design}

The Court agent assesses process compliance only, mirroring how UK courts
review procurement decisions: whether the authority followed a lawful,
rational and procedurally correct process, not whether the better bidder won
\cite{arrowsmith2014law}. Using a model to score another model's output is a
common evaluation pattern \cite{zheng2023judging}, and the same caution
applies, since the judge inherits the biases of the thing judged. Getting
the agent to behave took four revisions, each fixing a failure mode the
last one exposed --- from tracking dialogue concession rather than fact, to
inventing point systems on qualitative disputes, to the current conditional
design (Appendix~\ref{app:repro} tabulates all four).

V4 is derived from V3 through a single verified string substitution of one
named block, and the module raises an error at import if that anchor stops
matching. The two prompts therefore differ in exactly one variable, turning
the ablation in Section~\ref{sec:ablation} into a controlled comparison
rather than a before-and-after anecdote.

Given a round's transcript and scenario, the gating procedure checks
whether the scenario states a formula \emph{and} every input it needs: both
present triggers exact recomputation, a mismatch counting as manifest
error (Step 2A); neither present triggers qualitative rational-basis
review with no invented numbers (Step 2B); and formula-only --- the least
obvious branch, added after Woods v Milton Keynes exposed it --- falls
back to 2B rather than assuming or ``illustrating'' the missing inputs. A scenario stating a formula and a correction but not the original
inputs is not computable --- an earlier prompt let the model paper over
that gap by assuming the missing figures, producing arithmetic that looked
valid and was evidentially worthless. Step 1's gate now requires the
formula \emph{and} every input it needs, and bans writing down an unstated
number even hedged as illustrative: the hedge is not a safeguard, it is the
exact shape the failure took.

The same discipline applies to legal citation: the agent may only cite a
provision appearing in the scenario, in a party's statement that round, or
in its own system prompt's grounding material. A correct general point with
no citation is preferred over a specific-sounding invented one, since a
fabricated regulation number is the same category of error as a fabricated
sub-score. Section~\ref{sec:citation} measures whether that instruction
holds up.

\section{Orchestration}
\label{sec:orchestration}

Turn order runs through a LangGraph \texttt{StateGraph} rather than a
conversational free-for-all: \texttt{pre\_negotiation} sets out both
parties' interests, goals and BATNA, then \texttt{ca\_round} and
\texttt{bidder\_round} alternate with \texttt{court\_check} assessing
compliance each round, before \texttt{win\_statements} and a closing
\texttt{summary}. Everything hinges on \texttt{court\_check}: the only node
that can set \texttt{resolved} and route the loop back rather than forward
to the win statements, so the only escape from the loop other than
exhausting the round limit.

Resolution is decided by the Court's \texttt{recommended\_action}, drawn from
a closed vocabulary: continue negotiation, re-evaluation, no remedy with the
decision standing, or damages. Anything but ``continue'' resolves the
dispute; exhausting the round limit is recorded as deadlock. Constraining
this vocabulary is a legal requirement, not convenience --- Doc 3 is
explicit that award-stage disputes cannot be settled by splitting the
difference --- and it is enforced, not just documented: the downstream
module raises on any outcome outside the set, verified by injecting
``settle for 50\% of contract value'' during testing.

The orchestrator exposes both a synchronous \texttt{run()} for batch scripts
and a \texttt{stream()} generator for the API. One real bug: streamed
per-node updates are partial, and \texttt{bidder\_round} does not carry
\texttt{round\_number} the way \texttt{ca\_round} does, so a consumer must
read that field from merged state, not the raw update.

\section{From documents to scenarios}
\label{sec:extraction}

Real users upload PDFs and Word files, not Python objects. The extraction
agent turns raw text into a validated \texttt{DisputeScenario} in one model
call, with a repair retry on failure. Its prompt mirrors the Court agent's
discipline in both directions: do not invent a scoring formula the source
lacks, and do not summarise away one it has --- the second half was added
after a real case lost its published 60/40 weighting during extraction,
silently downgrading the downstream Court assessment from exact
verification to qualitative review.

What extraction itself does \emph{not} do is remove anything: no pass
strips the eventual court result from a source document, so a judgment's
own disposition can survive into the scenario the agents receive.
\S\ref{sec:agreement} shows how often that happened, and how it was fixed
rather than merely flagged.

Two guards were added after failures on documents I had not hand-picked: a
minimum source-length check raises before any model call when extraction
yields under 200 characters, because a scanned PDF with no text layer had
produced not an error but a fictional dispute; and an unstated contract
value must serialise as \texttt{null} rather than zero, since zero reads
downstream as a genuine \pounds0 contract.

\section{Structured output and compliance instrumentation}
\label{sec:compliance-method}

Every agent returns JSON validated against a Pydantic schema, with Ollama's
JSON mode constraining generation at the decoding step in the manner
Willard and Louf \cite{willard2023guided} formalise. Before this
instrumentation existed, a run where every response needed repair looked
identical to one where none did, so a counter layer was added around
parsing, recording brace-extraction fallbacks, field coercions, repetition
regenerations and outright parse failures. The headline metric is
\emph{structural compliance}: the proportion needing no repair, scoped per
response rather than summed over events, since one response can trigger
several coercions and naively summing them gives a rate that climbs past
one and goes negative. Repetition retries are counted but do not mark a
response dirty, since repetitiveness is a content problem, not a
structural one.

The same discipline extends to the
write-up:\label{sec:verification-discipline} no number, citation or case
fact enters this project without being checked against a primary source,
applied as strictly to my own prose as to model output. A contract value of
roughly \pounds1.5 billion attributed to one case proved, on checking
against the raw case text, to be \pounds112.1 million for the procurement at
issue --- the same class of error the Court agent is hardened against,
occurring in the human layer, and reported rather than fixed silently.

\section{Method selection, and what was rejected}
\label{sec:alternatives}

Table~\ref{tab:alternatives} records the methodologically significant
choices against the alternatives considered; two are worth commenting on
directly. Further plan-change decisions --- the frontend framework, the log
storage format, and the move from synthetic to real evaluation scenarios ---
are recorded in Appendix~\ref{app:repro}.

\begin{table}[htbp]
\centering
\footnotesize
\caption[Design alternatives]{Design decisions against the alternatives
considered. Rows marked \emph{plan change} diverge from the original plan.}
\label{tab:alternatives}
\begin{tabularx}{\textwidth}{@{}>{\raggedright\arraybackslash}p{2.9cm}>{\raggedright\arraybackslash}p{3.0cm}X@{}}
\toprule
\textbf{Decision} & \textbf{Alternative} & \textbf{Reasoning} \\
\midrule
Prompt engineering + context injection & Fine-tuning a domain model &
 No labelled corpus, and prompts stay auditable for fabrication testing. \\
Llama 3.1 8B \cite{grattafiori2024llama}, self-hosted \emph{(plan change)} & GPT-4o via API &
 Cost and data sensitivity; also the harder test, since discipline that
 holds on an 8B model isn't borrowed capability. \\
Single-pass extraction & Retrieval-augmented generation &
 RAG \cite{lewis2020rag} retrieves the right authority but does not stop
 invented facts about the case itself --- the actual failure mode here,
 scoped out deliberately. \\
LangGraph state machine & Free-form agent conversation &
 A dispute has a fixed procedural order, which an explicit graph encodes
 rather than constrains; it also enables the node-deletion ablation in
 \S\ref{sec:baselines}. \\
\bottomrule
\end{tabularx}
\end{table}

The move from synthetic to real scenarios was the most consequential
decision of this kind, since the model-authored synthetic pair made the
test circular and was deleted outright rather than kept with a caveat
(Appendix~\ref{app:further-decisions} gives the full reasoning). The plan
had also proposed ablations over temperature and BATNA values; those were
dropped in favour of the three baselines in Section~\ref{sec:baselines}, on
the judgement that ``does this architecture beat a single prompt'' is the
question an examiner asks first, not a temperature sweep.

\section{Delivery and the industrial deliverable}
\label{sec:risk-screen}

The simulator is exposed through the FastAPI backend and React client
above, with negotiations streamed live over server-sent events.
The component Fusion21 actually asked for is different: their framing was
prevention, since everything else here operates on a dispute already
raised. The pre-award challenge risk screen is a rule-based function with no
model call, taking the 16 profile fields encoded from Doc 1 and returning a
risk band, a normalised score, and a list of flags, each naming its
triggering field, Doc 1 category, severity, confidence tag, a rationale
paraphrased from Doc 1, and a concrete pre-award mitigation.

Three boundaries are deliberate. It predicts whether a challenge is likely
to be \emph{raised}, not whether the authority would lose one --- that
belongs to the Court agent, and blurring the two would smuggle a merits
judgment in through inputs that speak only to incentive to complain. Every
flag carries a confidence tag, since an authority knows its own
documentation quality but can only guess at a bidder's legal
representation. And the score is a triage ordering device, not a
probability: Doc 1 gives directional relationships, not weights, so the
severity weighting is a transparent first pass that says so rather than
claiming calibration it lacks.

